\subsection*{1.2. The geometry of Hilbert spaces}
\hfill 11

As a first consequence we obtain the \textbf{Cauchy-Schwarz-Bunjakowski}
inequality:

\noindent\textbf{Theorem 1.5 (Cauchy-Schwarz-Bunjakowski).} Let $\mathcal{H}_0$ be an inner product
space, then for every $f, g \in \mathcal{H}_0$ we have
\[
|\langle f,g \rangle| \leq \|f\|\|g\|
\tag{1.26}
\]
with equality if and only if $f$ and $g$ are parallel.

\noindent\textbf{Proof.} It suffices to prove the case $\|g\| = 1$. But then the claim follows
from $\|f\|^2 = |\langle g, f \rangle|^2 + \|f_{\perp}\|^2$.

Note that the Cauchy-Schwarz inequality implies that the scalar product
is continuous in both variables, that is, if $f_n \to f$ and $g_n \to g$ we have
$\langle f_n, g_n \rangle \to \langle f, g \rangle$.

As another consequence we infer that the map $\|\cdot\|$ is indeed a norm.
\[
\|f + g\|^2 = \|f\|^2 + \langle f, g \rangle + \langle g, f \rangle + \|g\|^2 \leq (\|f\| + \|g\|)^2.
\tag{1.27}
\]
But let us return to $\mathcal{C}(I)$. Can we find a scalar product which has the
maximum norm as associated norm? Unfortunately the answer is no! The
reason is that the maximum norm does not satisfy the parallelogram law
(Problem 1.7).

\noindent\textbf{Theorem 1.6 (Jordan-von Neumann).} A norm is associated with a scalar
product if and only if the \textbf{parallelogram law}
\[
\|f + g\|^2 + \|f - g\|^2 = 2\|f\|^2 + 2\|g\|^2
\tag{1.28}
\]
holds.

In this case the scalar product can be recovered from its norm by virtue
of the \textbf{polarization identity}
\[
\langle f,g \rangle = \frac{1}{4} \left( \|f + g\|^2 - \|f – g\|^2 + i\|f + ig\|^2 – i\|f + ig\|^2 \right).
\tag{1.29}
\]

\noindent\textbf{Proof.} If an inner product space is given, verification of the parallelogram
law and the polarization identity is straight forward (Problem 1.6).
To show the converse, we define
\[
s(f,g) = \frac{1}{4} \left( \|f + g\|^2 - \|f – g\|^2 + i\|f + ig\|^2 – i\| f + ig\|^2 \right).
\tag{1.30}
\]
Then $s(f, f) = \|f\|^2$ and $s(f,g) = s(g, f)^*$ are straightforward to check.
Moreover, another straightforward computation using the parallelogram law
shows
\[
s(f,g) + s(f,h) = 2s\left(f, \frac{g+h}{2}\right).
\tag{1.31}
\]